\subsection{Criterios de lisitud y no ramificación por las fibras}
\label{subsection:IV.17.8}
\label{IV.17.8}

\begin{proposition}[17.8.1]
\label{IV.17.8.1}
Sean $g:Y\to S$, $h:X\to S$ dos morfismos localmente de presentación finita. Para que un $S$-morfismo $f:X\to Y$ sea no ramificado, es necesario y suficiente que, para todo $s\in S$, el morfismo $f_s:h^{-1}(s)\to g^{-1}(s)$ deducido de $f$ por el cambio de base $\Spec(k(s))\to S$ sea no ramificado.
\end{proposition}

Esto resulta inmediatamente del hecho de que, para un morfismo localmente de presentación finita, ser no ramificado es una propiedad de las \emph{fibras} de dicho morfismo (17.4.1, $d$)), y de que, para todo $y\in Y$, se tiene $f^{-1}(y)=f_s^{-1}(y)$ si $s=g(y)$.

\begin{proposition}[17.8.2]
\label{IV.17.8.2}
Sean $g:Y\to S$, $h:X\to S$ dos morfismos localmente de presentación finita; se supone además que $h$ es plano. Para que un $S$-morfismo $f:X\to Y$ sea liso (resp. étale), es necesario y suficiente que, para todo $s\in S$, el morfismo $f_s:h^{-1}(s)\to g^{-1}(s)$ deducido de $f$ por el cambio de base $\Spec(k(s))\to S$ sea liso (resp. étale). Cuando esto ocurre, el morfismo $g$ es plano en los puntos de $f(X)$.
\end{proposition}

En efecto, se sabe (11.3.10) que, para que $f$ sea plano, es necesario y suficiente que $f_s$ lo sea para todo $s\in S$, y que entonces $g$ es plano en los puntos de $f(X)$. Pero, para un morfismo plano localmente de presentación finita, ser liso es una propiedad de las fibras de dicho morfismo (17.5.1, $b$)), y para todo $y\in Y$, se tiene $f^{-1}(y)=f_s^{-1}(y)$ si $s=g(y)$.

\begin{remark}[17.8.3]
\label{IV.17.8.3}
Las demostraciones anteriores muestran (teniendo en cuenta (11.3.10)) que, bajo las mismas hipótesis anteriores sobre $g$ y $h$, para que $f$ sea no ramificado (resp. liso, resp. étale) \emph{en un punto} $x\in X$, basta que, si se pone $s=h(x)$, $f_s$ sea no ramificado (resp. liso, resp. étale) en el punto $x$.
\end{remark}
